\NPHardness*
\begin{proof}

\end{proof}


\NPHardnessOfClash*
\begin{proof}
Let us fix a parameter $n$. Before the reduction, we construct a gadget that distinguishes all words of length at most $n$ unless they consist one of the designated letters. 
Let $Y \subseteq \Sigma$ and $x \in \Sigma \setminus Y$. Consider the automaton $\autB_{Y,x}$ depicted below:

\begin{figure}[h!]
\centering
\begin{tikzpicture}[
    shorten >=1pt, 
    node distance=2.5cm, 
    on grid, 
    auto,
    >={Stealth[bend]}
]

  % State Definitions
  \node[state] (qn) {$q_n$};
   \node[draw=none] (dots) [right=of qn] {$\dots$};
  \node[state] (q2) [right=of dots] {$q_2$};
  \node[state] (q1) [right=of q2] {$q_1$};
\node[state] (qF) [right=of q1] {$q_F$};

\node[state] (sink) [below=of q2] {$s$};

  % Transition Edges
  \path[->]
    (qn) edge node [above] {$\Sigma\setminus Y$} (dots)
    (dots) edge node [above] {$\Sigma\setminus Y$} (q2)
    (q2)  edge node [above] {$\Sigma\setminus Y$} (q1)
    % The top arc from qn back to q1
    (q1)  edge [bend right=45] node [above] {$\Sigma\setminus (Y \cup \set{x})$} (qn)
    % The outgoing edge xi from q1
%    (q1)  edge [nodes={at={(1, -0.5)}}] node {$x$} ++(1.5, 0);
    (q1)  edge node [above] {$x$} (qF)
    (q1) edge node [left] {$Y$} (sink)
    (q2) edge node [left] {$Y$} (sink)
    (qn) edge node [above] {$Y$} (sink)
    (qF) edge node [left] {$Y$} (sink)
    ;

 \draw[loop,->] (qF)  to node[above] {$\Sigma \setminus Y$} (qF);
 \draw[loop right,->] (sink) to node[right] {$\Sigma$} (sink);

\end{tikzpicture}
%\label{The gadget \DFA $\autB_{Y,x}$, which parially distingushes words of length at most $n$, which do not contain any letter from $Y$ }
\end{figure}

Observe that for all words $w_1, w_2$ of length bounded by $n$, if $f^{\autB_{Y,x}}_{w_1} = f^{\autB_{Y,x}}_{w_2}$, then one of the following holds:
\begin{itemize}
    \item each of the words $w_1, w_2$ contains at least one letter from $Y$; in that case $f_{w_1}, f_{w_2}$ are constant functions equal 
    $s$ for every argument,
    \item the sets positions at which $x$ occurs in $w_1$ and resp. $w_2$ are equal; 
    indeed, for any $w \in (\Sigma \setminus Y)^{\leq n}$ we have
    $f_{w}(q_i) = q_F$ if and only if $w[i] = x$.
\end{itemize}
Now, consider $\autB_Y$ being the disjoint union of automata $\autB_{Y,x}$ over all $x \in \Sigma\setminus Y$.

Observe that for all words $w_1, w_2$ of length bounded by $n$, if $f^{\autB_{Y}}_{w_1} = f^{\autB_{Y}}_{w_2}$, then one of the following holds:
\begin{itemize}
    \item each of the words $w_1, w_2$ contains at least one letter from $Y$; in that case $f_{w_1}, f_{w_2}$ are constant functions equal 
    $s$ for every argument,
    \item words $w_1$ and $w_2$ are equal.
\end{itemize}

Having the automaton $\autB_Y$ we reduce the 3-SAT problem to the minimal collision problem.
Let $\varphi$ be a 3-CNF formula over variables $x_1, \ldots, x_n$.
Consider $\Sigma = \{x_1, \bar{x_1}, \ldots, x_n, \bar{x_n}, a, b \}$.
With each clause $c_j$ we associate the subset $C_j$ of $\Sigma$ consisting of the literals from $c_j$.
Furthermore, for every variable $x_i$ we define the set $X_i = \set{x_i, \bar{x_i}}$.
Now, consider a \DFA $\aut_0$ being the disjoint union of automata $\autB_{C_1}, \ldots, \autB_{C_m}$ over all clauses of $\varphi$ and
automata $\autB_{X_1}, \ldots, \autB_{X_n}$ over all variables of $\varphi$. 
We extend $\aut_0$ to $\aut$ with additional states such that over these states letters from $\Sigma \setminus \set{a,b}$ are self loops, and
letters $a,b$ ensure that every component of the disjoint union is reachable. 
Let $k$ in the instance of the minimal collision problem be equal $n$, the number of variables.

Now, consider a pair of different words $w_1, w_2$ such that $f^{\aut}_{w_1} = f^{\aut}_{w_2}$ and $|w_1|,|w_2| \leq n$.
Then, for every automaton $\autB_{X_1}, \ldots, \autB_{X_n}$, since $w_1 \neq w_2$, each word contains at least one letter from each set $X_i$.
There are $n$ such sets, which are disjoint, and hence $w_1$ and $w_2$ contain exactly one letter from each of the sets $X_i$, i.e., 
$w_1$ and $w_2$ represent consistent variable assignments (they either contain $x_i$ or $\bar{x_i}$).
Furthermore, for every clause $c_j$ each word $w_1, w_2$ contains at least one letter corresponding to some literal from $c_j$.
Therefore, $w_1$ and $w_2$ represent (possibly different) satisfying assignments for $\varphi$.

Conversely, assume that $\varphi$ has a satisfying assignment. 
Consider two words $w_1, w_2$ representing this variable assignment with a different order or literals.
Then, for every component $\autB_{Y,x}$ of $\aut_0$, the functions $f_{w_1}^{\autB_{Y,x}}, f_{w_2}^{\autB_{Y,x}}$ are constant functions equal $s$.
Furthermore, on states in $\aut$, which are not in $\aut_0$, both functions $f_{w_1}^{\aut}, f_{w_2}^{\aut}$ coincide with the identity.
Consequently, we have
(a)~$w_1 \neq w_2$, (b)~$|w_1| = |w_2| = n$, and
(c)~$f_{w_1}^{\aut_0}, f_{w_2}^{\aut_0}$, i.e., 
the constructed instance of the minimal collision problem is positive.
\end{proof}